% ---
% id: foundations/cor:rank-over-factor-enclosure
% kind: cor
% title: "Rank over a factor enclosure"
% summary: >-
%   A factor enclosure with no column holds only the zero matrix. One whose
%   Gram matrix L*L has a positive floor over its box fixes, at every member,
%   the rank, its number of columns, and the support, the span of that
%   member's columns.
% prominence: technical
% status: implemented
% uses: [def:factor-enclosure, cor:least-eigenvalue-over-enclosure]
% mentions: []
% assumes: []
% refs: []
% provenance: newly-derived
% implements: certiqs_key.foundations.reference:fixes_rank
% lean:
% ---

Corollary~\ref{cor:least-eigenvalue-over-enclosure}\,(iv) fixes the rank of
the matrices an entrywise enclosure contains where they are positive definite
on the rows it does not fix at zero. A \tref{factor enclosure} fixes it where
the columns of its factors are linearly independent throughout the box, which
a positive floor under their Gram matrix shows: every member then has as many
independent columns as the triple has, and its support is the span of its own
columns. The support moves with the member, its dimension does not, and each
member has a factor form whose factor lies in a box and whose core is the same
exact matrix for all. Part~(iii) gives one floor that is a comparison of
rationals.

\begin{corollary}\label{cor:rank-over-factor-enclosure}
Let $(L_c,t,W)$ be a triple of a \tref{matrix box} $(L_c,t)$ of complex
$d\times k$ matrices and a diagonal $k\times k$ matrix $W$ with positive
rational diagonal entries, as in Definition~\ref{def:factor-enclosure}, and
call a matrix $L$ that the box contains a \emph{factor} and $LWL^{\dagger}$ a
\emph{member}.
\begin{enumerate}[label=(\roman*),leftmargin=2.4em]
\item \emph{(No column.)} If $k=0$, every member is the zero $d\times d$
  matrix.
\item \emph{(A Gram floor fixes the rank.)} Let $k\ge1$, let $G$ be a
  Hermitian $k\times k$ matrix and $\rho\ge0$ with
  $\|L^{\dagger}L-G\|_{\mathrm F}\le\rho$ for every factor $L$, the Frobenius
  norm of Corollary~\ref{cor:least-eigenvalue-over-enclosure}, and let
  $\ell\le\lmin(G)$ with $\ell-\rho>0$. Then for every factor $L$,
  $L^{\dagger}L\succ0$, so that $k\le d$ and $\operatorname{rank}L=k$; and the
  member $X\coloneqq LWL^{\dagger}$ satisfies $X=E^{\dagger}WE$ with
  $E\coloneqq L^{\dagger}$ a $k\times d$ matrix of rank $k$, $X\succeq0$,
  $\operatorname{ran}X=\operatorname{ran}L$, the span of the columns of $L$,
  and $\operatorname{rank}X=k$; in particular $X\neq0$.
\item \emph{(One such floor.)} Let $k\ge1$, put
  $a_{im}\coloneqq|\operatorname{Re}(L_c)_{im}|+|\operatorname{Im}(L_c)_{im}|$
  and
  \begin{equation}\label{eq:gram-radius}
    s_{ml}\;\coloneqq\;\sum_{i=1}^{d}\big(a_{im}t_{il}+t_{im}a_{il}
    +t_{im}t_{il}\big)\qquad(1\le m,l\le k),
  \end{equation}
  rationals. Then $G\coloneqq L_c^{\dagger}L_c$ is Hermitian with entries in
  $\mathbb{Q}(i)$, and every rational $\rho\ge0$ with
  $\rho^{2}\ge\sum_{m,l}s_{ml}^{2}$ satisfies
  $\|L^{\dagger}L-G\|_{\mathrm F}\le\rho$ for every factor $L$; whether $\rho$
  qualifies, and whether $\ell-\rho>0$ for a rational $\ell$, is a comparison
  of rationals.
\end{enumerate}
\end{corollary}

\begin{proof}
\emph{(i).} For $k=0$ a factor $L$ is $d\times0$ and $W$ is $0\times0$, so
every entry $(LWL^{\dagger})_{il}=\sum_{m=1}^{k}L_{im}W_{mm}\overline{L_{lm}}$
is an empty sum, $0$.

\emph{(ii).} Let $L$ be a factor. $L^{\dagger}L$ is a Hermitian $k\times k$
matrix, and Corollary~\ref{cor:least-eigenvalue-over-enclosure}\,(i) with
$R\coloneqq L^{\dagger}L$ and $M\coloneqq G$ gives
$\lmin(L^{\dagger}L)\ge\lmin(G)-\rho\ge\ell-\rho>0$. So for $y\ne0$,
$\|Ly\|^{2}=y^{\dagger}L^{\dagger}Ly>0$: $L$ is injective, its rank is $k$,
and $k\le d$, the rank of a $d\times k$ matrix being at most $d$. The adjoint
$E=L^{\dagger}$ has the same rank $k$, and $E^{\dagger}WE=LWL^{\dagger}=X$
with $W$ positive definite, its diagonal entries being positive.
Corollary~\ref{cor:least-eigenvalue-over-enclosure}\,(iii), with $E$ and
$\widetilde X\coloneqq W$, gives $X\succeq0$,
$\operatorname{ran}X=\operatorname{ran}E^{\dagger}=\operatorname{ran}L$ and
$\operatorname{rank}X=k\ge1$, so $X\ne0$.

\emph{(iii).} The entries $(L_c^{\dagger}L_c)_{ml}=\sum_i
\overline{(L_c)_{im}}(L_c)_{il}$ are sums of products of Gaussian rationals,
so they lie in $\mathbb{Q}(i)$, and
$\overline{(L_c^{\dagger}L_c)_{lm}}=(L_c^{\dagger}L_c)_{ml}$. Let $L$ be a
factor and $\Delta\coloneqq L-L_c$, so $|\Delta_{im}|\le t_{im}$. Then
\begin{equation*}
  (L^{\dagger}L-L_c^{\dagger}L_c)_{ml}
  =\sum_{i=1}^{d}\Big(\overline{(L_c)_{im}}\,\Delta_{il}
  +\overline{\Delta_{im}}\,(L_c)_{il}+\overline{\Delta_{im}}\,\Delta_{il}\Big),
\end{equation*}
and $|(L_c)_{im}|\le a_{im}$, since $|x+iy|\le|x|+|y|$ for real $x,y$; by the
triangle inequality the modulus of the left-hand side is at most $s_{ml}$.
Hence $\|L^{\dagger}L-G\|_{\mathrm F}^{2}=\sum_{m,l}|(L^{\dagger}L-G)_{ml}|^{2}
\le\sum_{m,l}s_{ml}^{2}\le\rho^{2}$. The $a_{im}$ and $s_{ml}$ are sums and
products of rationals, and so is $\sum_{m,l}s_{ml}^{2}$.
\end{proof}
